\section{Resultados e Discussão}

\subsection{Métricas de Avaliação}
Para analisar o desempenho dos classificadores, foram utilizadas métricas fundamentais de aprendizado de máquina: \textit{Precision}, \textit{Recall}, \textit{F1-Score} e \textit{Support}.

A \textit{Precision} avalia a porcentagem de instancias dadas como positivas que eram realmente positivas, ajudando a identificar a quantidade adivinhações falsas positivas. Em contrapartida, o \textit{Recall} mede a capacidade do modelo de identificar corretamente todas as instancias positivas presentes no dataset, servindo para medir a proporção de adivinhações falsas negativas.

O \textit{F1-Score} consiste na média harmônica entre a precisão e o recall, sendo particularmente útil quando se busca um equilíbrio entre essas duas métricas. Por fim, o \textit{Support} refere-se ao número de ocorrências de cada instancia no dataset tanto negativas(Normal) quanto negativas(Anomala)

A métrica de \textit{acurácia} não é ideal para essa analise tendo em vista que a maioria dos acessos de rede no mundo real são normais e raramente anômalos, então é muito mais interessante analisar falsos negativos e positivos do que o acerto geral.

\subsection{Resultados com K-Nearest Neighbors (KNN)}

O K-Nearest Neighbors (KNN) foi avaliado com a configuração otimizada de $K=3$ e métrica de distância 
Manhattan. No cenário utilizando o conjunto completo de atributos, o modelo alcançou uma acurácia global 
de 0.9943. O classificador demonstrou alta capacidade de generalização e separação, com F1-Score de 0.99 
para ambas as classes, minimizando de forma eficaz os falsos positivos e falsos negativos, conforme 
evidenciado na Tabela~\ref{tab:knn_full}.

\begin{table}[H]
\centering
\caption{Métricas de desempenho do KNN (dataset completo)}
\label{tab:knn_full}
\begin{tabular}{lcc}
\hline
\textbf{Métrica} & \textbf{Classe Normal} & \textbf{Classe Anomalia} \\
\hline
Precision & 0.99 & 1.00 \\
Recall    & 1.00 & 0.99 \\
F1-Score  & 0.99 & 0.99 \\
Support   & 4035 & 3523 \\
\hline
\end{tabular}
\end{table}

Apesar do excelente poder discriminativo, a avaliação revelou um fator limitante crítico inerente à arquitetura 
de aprendizado tardio (\textit{lazy learning}) do KNN: o elevado tempo de inferência. O modelo demandou 
aproximadamente 3,36 segundos para classificar as instâncias de teste. Trata-se de um custo computacional 
significativamente superior ao dos modelos paramétricos avaliados, o que pode representar um gargalo severo em 
aplicações de detecção de intrusão em tempo real.

Com o objetivo de avaliar a robustez do modelo frente a conjuntos de dados incompletos, o KNN foi reavaliado após 
a exclusão da feature \textit{src\_bytes}. Diferentemente da queda de desempenho observada em classificadores 
baseados em árvores de decisão, o KNN apresentou uma resiliência absoluta. Conforme apresentado na 
Tabela~\ref{tab:knn_reduced}, todas as métricas permaneceram rigorosamente inalteradas, mantendo a acurácia em 0.9943.

\begin{table}[H]
\centering
\caption{Métricas de desempenho do KNN (dataset reduzido, sem \textit{src\_bytes})}
\label{tab:knn_reduced}
\begin{tabular}{lcc}
\hline
\textbf{Métrica} & \textbf{Classe Normal} & \textbf{Classe Anomalia} \\
\hline
Precision & 0.99 & 1.00 \\
Recall    & 1.00 & 0.99 \\
F1-Score  & 0.99 & 0.99 \\
Support   & 4035 & 3523 \\
\hline
\end{tabular}
\end{table}

Como as inferências foram realizadas sobre dados de teste, descarta-se a hipótese de \textit{overfitting}. Este 
comportamento altamente estável é explicado por propriedades estruturais do modelo operando neste espaço de dados: 
(1) a diluição vetorial em um espaço de alta dimensionalidade (118 dimensões resultantes da codificação binária), onde 
a exclusão de um único eixo afeta de forma negligenciável o cálculo geral de similaridade; (2) a redundância informacional 
e forte correlação inerente às assinaturas de ataques de rede; e (3) a robustez da inferência baseada em distância 
geométrica global, que, ao contrário das regras de partição locais, não colapsa com a ausência de um único atributo de 
alto ganho de informação.

\subsection{Resultados com Isolation Forest}

O Isolation Forest se mostrou ineficiente quando comparado aos outros
modelos utilizados durante a pesquisa, apresentando a menor acurácia (0.748).
Entretanto, como falado anteriormente, a acurácia não é a métrica mais confiável em conjuntos de dados
desbalanceados, o que é comum em problemas de detecção de anomalias, pois pode ser
inflacionada pela capacidade do modelo em classificar corretamente a classe
majoritária (tráfego normal). Dessa forma, uma análise mais aprofundada torna-se
necessária.

\begin{table}[ht]
\centering
\caption{Métricas de desempenho do Isolation Forest}
\label{tab:if_metrics}
\begin{tabular}{lcc}
\hline
\textbf{Métrica} & \textbf{Classe Normal} & \textbf{Classe Anomalia} \\
\hline
Precision & 1.00 & 0.65 \\
Recall    & 0.53 & 1.00 \\
F1-Score  & 0.69 & 0.79 \\
Support   & 2690 & 2349 \\
\hline
\end{tabular}
\end{table}

Como pode ser observado na Tabela~\ref{tab:if_metrics}, todos os eventos classificados
como normais eram, de fato, normais. A principal dificuldade do modelo esteve na
identificação correta das anomalias, apresentando precisão de 65\% (0.65) ao rotular
eventos como intrusão. Além disso, considerando todos os eventos analisados, o modelo
foi capaz de identificar corretamente apenas 53\% (0.53) dos eventos normais, o que
indica a geração de um elevado número de falsos positivos, isto é, eventos normais
classificados como ataques.

Por outro lado, o Isolation Forest não apresentou o problema de gerar falsos negativos,
uma vez que ataques não foram classificados como eventos normais. Em síntese, o modelo
apresenta dificuldades em verificar se um evento é realmente normal, gerando falsos
positivos que podem ocasionar retrabalho para equipes de segurança, as quais precisam
analisar manualmente eventos normais identificados como anômalos.

\subsection{Resultados com Random Forest}

O Random Forest apresentou o melhor desempenho entre todos os modelos avaliados quando
utilizado com o conjunto completo de atributos, alcançando uma acurácia de 0.9975. O
modelo demonstrou excelente capacidade de generalização, classificando corretamente
tanto o tráfego normal quanto as intrusões, o que evidencia seu elevado poder
discriminativo em um cenário de detecção de anomalias em redes.

Conforme apresentado na Tabela~\ref{tab:rf_metrics_full}, os valores de Precision,
Recall e F1-Score foram iguais a 1.00 para ambas as classes, indicando a ausência
praticamente total de falsos positivos e falsos negativos. Esses resultados reforçam a
robustez do Random Forest e justificam sua utilização como modelo de referência para a
análise de importância dos atributos. Nesse contexto, a feature \textit{src\_bytes} foi
identificada como a mais influente no modelo, motivando sua remoção em uma segunda etapa
experimental.

Com o objetivo de avaliar a robustez do Random Forest frente a conjuntos de dados
incompletos, o modelo foi reavaliado após a exclusão da feature \textit{src\_bytes}.
Nesse cenário, observou-se uma queda significativa no desempenho, com a acurácia
reduzida para 0.9630, conforme apresentado na Tabela~\ref{tab:rf_metrics_reduced}.
Embora o modelo ainda apresente bons resultados globais, nota-se uma redução no Recall
da classe normal (0.93), indicando um aumento na taxa de falsos positivos, ou seja,
eventos normais classificados incorretamente como anomalias.

Esse comportamento evidencia a forte dependência do Random Forest em relação a
atributos altamente informativos e reforça a importância da qualidade e da completude
dos dados em aplicações de detecção de intrusões em redes. Além disso, os resultados
obtidos corroboram a motivação deste trabalho em investigar como diferentes modelos de
aprendizado de máquina reagem à ausência de atributos relevantes no conjunto de dados.

\begin{table}[ht]
\centering
\caption{Métricas de desempenho do Random Forest (dataset completo)}
\label{tab:rf_metrics_full}
\begin{tabular}{lcc}
\hline
\textbf{Métrica} & \textbf{Classe Normal} & \textbf{Classe Anomalia} \\
\hline
Precision & 1.00 & 1.00 \\
Recall    & 1.00 & 1.00 \\
F1-Score  & 1.00 & 1.00 \\
Support   & 3523 & 4035 \\
\hline
\end{tabular}
\end{table}


\begin{table}[ht]
\centering
\caption{Métricas de desempenho do Random Forest (dataset reduzido, sem \textit{src\_bytes})}
\label{tab:rf_metrics_reduced}
\begin{tabular}{lcc}
\hline
\textbf{Métrica} & \textbf{Classe Normal} & \textbf{Classe Anomalia} \\
\hline
Precision & 0.99 & 0.94 \\
Recall    & 0.93 & 0.99 \\
F1-Score  & 0.96 & 0.97 \\
Support   & 3523 & 4035 \\
\hline
\end{tabular}
\end{table}


\subsection{Resultados com MLP}

O Multilayer Perceptron foi avaliado em dois cenários distintos: um utilizando todas as
features disponíveis no conjunto de dados e outro com a remoção da feature
\textit{src\_bytes}, identificada como a mais influente no modelo Random Forest, que
apresentou o melhor desempenho geral. O objetivo dessa análise foi avaliar o impacto da
ausência de uma feature relevante no desempenho do MLP, simulando cenários em que o
conjunto de dados pode estar incompleto.

No cenário completo, o MLP obteve uma acurácia de 0.9960, apresentando desempenho
equilibrado entre as classes. A classe Anomalia apresentou precisão de 0.9969 e recall
de 0.9946, enquanto a classe Normal obteve precisão de 0.9953 e recall de 0.9973,
resultando em F1-Scores superiores a 0.995 para ambas as classes, conforme apresentado
na Tabela~\ref{tab:mlp_full}.

\begin{table}[ht]
\centering
\caption{Métricas de desempenho do MLP com todas as features}
\label{tab:mlp_full}
\begin{tabular}{lcc}
\hline
\textbf{Métrica} & \textbf{Classe Anomalia} & \textbf{Classe Normal} \\
\hline
Precision & 0.9969 & 0.9953 \\
Recall    & 0.9946 & 0.9973 \\
F1-Score  & 0.9957 & 0.9963 \\
Support   & 3523   & 4035   \\
\hline
\end{tabular}
\end{table}

Após a remoção da feature \textit{src\_bytes}, o MLP apresentou uma leve redução de
desempenho, com acurácia de 0.9952. Observa-se que, embora o recall da classe Anomalia
tenha aumentado para 0.9969, houve redução na precisão dessa classe, passando para
0.9929. Para a classe Normal, verificou-se comportamento inverso, com aumento da
precisão e redução do recall. Ainda assim, os valores de F1-Score permaneceram elevados
para ambas as classes, indicando que o modelo manteve desempenho consistente mesmo com a
remoção de uma feature altamente relevante em outro classificador, conforme mostrado na
Tabela~\ref{tab:mlp_reduced}.

\begin{table}[ht]
\centering
\caption{Métricas de desempenho do MLP sem a feature \textit{src\_bytes}}
\label{tab:mlp_reduced}
\begin{tabular}{lcc}
\hline
\textbf{Métrica} & \textbf{Classe Anomalia} & \textbf{Classe Normal} \\
\hline
Precision & 0.9929 & 0.9973 \\
Recall    & 0.9969 & 0.9938 \\
F1-Score  & 0.9949 & 0.9955 \\
Support   & 3523   & 4035   \\
\hline
\end{tabular}
\end{table}

De forma geral, os resultados indicam que o MLP é relativamente robusto à remoção de uma
feature altamente influente em outro modelo, apresentando apenas pequenas variações nas
métricas de desempenho. Esse comportamento reforça a capacidade do MLP de aprender
representações distribuídas dos dados, reduzindo sua dependência de atributos
individuais.

